\begin{frame}{곱적: 소실·폭발의 씨앗}
  이 식이 말해주는 것: \textbf{먼 시점과 가까운 시점 사이를
  건너는 그래디언트는 $W_{hh}$를 $T-1$번 곱하는 곱적}이 된다.
  \vskip0.6em
  \begin{columns}[T]
    \begin{column}{0.5\textwidth}
      \begin{block}{$W_{hh}$ 고유값 $< 1$}
        곱이 \textbf{지수적으로 0}으로 --- ``그래디언트 \textbf{소실}''
      \end{block}
      \vskip0.5em
      \begin{block}{$W_{hh}$ 고유값 $> 1$}
        곱이 \textbf{지수적으로 $\infty$}로 발산 --- ``그래디언트
        \textbf{폭발}''
      \end{block}
    \end{column}
    \begin{column}{0.5\textwidth}
      11.3절에서 ``그래디언트 소실·폭발''이라는 이름으로 다시 만나는
      바로 그 문제 (Pascanu et al., 2013).
      \vskip0.5em
      \kq{이번 절: 그 \textbf{양}보다 먼저, 이 곱적 구조를
      \textbf{정확히} 따라가며 각 파라미터의 그래디언트가 어떻게
      \textbf{누적}되는지 유도한다.}
    \end{column}
  \end{columns}
\end{frame}
